% claim-ref: aquarium-hornwort#aquarium-gentle-flow-intake-clearance
% claim-ref: aquarium-hornwort#failure-cue-breakage
% claim-ref: aquarium-hornwort#failure-cue-heavy-needle-shedding
% claim-ref: aquarium-hornwort#failure-cue-no-growth
% claim-ref: aquarium-hornwort#fragment-cut-length
% claim-ref: aquarium-hornwort#fragment-establishment-interval
% claim-ref: aquarium-hornwort#fragment-parent-readiness
% claim-ref: aquarium-hornwort#root-initiation
% claim-ref: aquarium-hornwort#step-1-prepare
% claim-ref: aquarium-hornwort#step-2-select
% claim-ref: aquarium-hornwort#step-3-trim
% claim-ref: aquarium-hornwort#step-4-no-callus
% claim-ref: aquarium-hornwort#step-5-place-floating
% claim-ref: aquarium-hornwort#step-5-place-supported
% claim-ref: aquarium-hornwort#step-6-maintain
% claim-ref: aquarium-hornwort#step-7-inspect
% claim-ref: aquarium-hornwort#step-8-maintain-dispose
% claim-ref: aquarium-hornwort#success-established-growth
% claim-ref: aquarium-hornwort#success-new-growth
% claim-ref: aquarium-hornwort#terrestrial-callusing
% claim-ref: aquarium-hornwort#terrestrial-rooting-depth
\CompanionHeader{PROPAGATION}{Aquarium hornwort}{Provisional Ceratophyllum demersum}{Aquarium; species and trade form unverified. Tank and livestock needs govern all care.}
\Context{Submerged fragment / floating primary. General guidance; observe parent readiness before acting.}
\begin{minipage}[t]{.625\linewidth}\vspace{0pt}
\Step{1}{Prepare \& contain}{Prepare clean aquarium scissors and a container of tank water. Keep all material wet and contained while working; capture unintended fragments. \textit{Proposed starting point.} \textbf{[HOR-NYDEC; HOR-WA]}}
\Step{2}{Select viable material}{Choose an intact, actively growing leafy tip or side shoot already submerged. Reject crushed or disintegrating pieces. No supported minimum cutting length is established. \textit{Proposed starting point.} \textbf{[HOR-DEN; HOR-WA]}}
\Step{3}{Separate gently}{Make one gentle clean trim to separate the tip or side shoot. Collect every loose piece; breakage alone is not evidence of successful regrowth. \textit{Proposed starting point.} \textbf{[HOR-DEN; HOR-WA]}}
\Step{4}{Keep submerged}{Move directly between aquarium water and the aquatic placement. Do not dry, callus or apply rooting hormone; this plant is rootless. \textit{Proposed starting point.} \textbf{[HOR-WA]}}
\Step{5}{Float with clearance}{Float freely below the surface, keeping needles clear of the filter intake and leaving room to observe the tip. Avoid crushing it against equipment. \textit{Proposed starting point.} \textbf{[HOR-AQCOOP; HOR-WA]}}
\Step{6}{Maintain recorded conditions}{Keep conditions stable and livestock-compatible. Record temperature, pH/hardness, light schedule, flow, nutrients and CO2; do not change unknown conditions solely for the fragment. \textit{Proposed starting point.} \textbf{[HOR-DEN; HOR-TROP; HOR-WA]}}
\Step{7}{Inspect actual regrowth}{Look for new tip extension or a new side shoot. Establishment means the fragment stays intact and continues growing. There is no defensible universal establishment interval. \textit{Proposed starting point.} \textbf{[HOR-DEN; HOR-WA]}}
\Step{8}{Trim \& dispose responsibly}{Trim crowding, capture every piece and seal unwanted material in a bag for household trash. Never release or compost aquarium fragments. Continue the maintenance log. \textit{Proposed starting point.} \textbf{[HOR-DEN; HOR-NYDEC]}}
\end{minipage}\hfill\begin{minipage}[t]{.345\linewidth}\vspace{0pt}
\Note{Milestones, not deadlines}{Cut/break: starting event only. Regrowth: new tip or side shoot. Established: continued intact growth. Root initiation and terrestrial callusing: N/A. \textit{Proposed starting point.} \textbf{[HOR-DEN; HOR-WA]}}
\Note{Supported placement}{The alternate is bottom support without calling the stem rooted. No terrestrial burial depth applies; floating remains the primary method. \textbf{[HOR-TROP]}}
\Note{Check 1}{Repeated breakage: check handling, crushing and intake/current exposure. \textit{Proposed starting point.} \textbf{[HOR-WA]}}
\Note{Check 2}{No regrowth: verify identity and record tank conditions before one change. \textit{Proposed starting point.} \textbf{[HOR-DEN; HOR-TROP; HOR-WA]}}
\Note{Check 3}{Heavy shedding: check parameter shifts, flow/intake, shaded bases, nutrients and liquid carbon; vacuum needles. \textit{Proposed starting point.} \textbf{[HOR-AQCOOP]}}
\end{minipage}\par\vspace{4pt}
\Panel{Ready for normal care}{Continued new growth on an intact submerged fragment under recorded conditions. No terrestrial transplant, rooting depth or soil mixture applies.\par Start date: \hrulefill\quad Method: \hrulefill\par First successful growth / conditions: \hrulefill}
\Sources{Evidence: HOR-AQCOOP Aquarium Co-Op; HOR-DEN Dennerle Plants GmbH; HOR-NYDEC New York State Department of Environmental Conservation; HOR-TROP Tropica Aquarium Plants; HOR-WA Washington State Department of Ecology. Full titles, URLs and claim scopes: entry sources.yaml and companion worksheets. Conversions are rounded; source units remain authoritative.}{aquarium-hornwort / propagation}
